\documentclass[12pt]{amsart}
%prepared in AMSLaTeX, under LaTeX2e
\addtolength{\oddsidemargin}{-.65in} 
\addtolength{\evensidemargin}{-.65in}
\addtolength{\topmargin}{-.4in}
\addtolength{\textwidth}{1.3in}
\addtolength{\textheight}{.8in}

\renewcommand{\baselinestretch}{1.05}

\usepackage{verbatim} % for "comment" environment

\usepackage{palatino}

\newtheorem*{thm}{Theorem}
\newtheorem*{defn}{Definition}
\newtheorem*{example}{Example}
\newtheorem*{problem}{Problem}
\newtheorem*{remark}{Remark}

\usepackage{fancyvrb,xspace,dsfont}

\usepackage[final]{graphicx}

% macros
\usepackage{amssymb}
\usepackage[dvipsnames]{xcolor}

\usepackage{hyperref}
\hypersetup{pdfauthor={Ed Bueler},
            pdfcreator={pdflatex},
            colorlinks=true,
            citecolor=blue,
            linkcolor=red,
            urlcolor=blue,
            }

\newcommand{\br}{\mathbf{r}}
\newcommand{\bv}{\mathbf{v}}
\newcommand{\bx}{\mathbf{x}}
\newcommand{\by}{\mathbf{y}}

\newcommand{\cB}{\mathcal{B}}
\newcommand{\cD}{\mathcal{D}}
\newcommand{\cE}{\mathcal{E}}
\newcommand{\cF}{\mathcal{F}}
\newcommand{\cH}{\mathcal{H}}
\newcommand{\cL}{\mathcal{L}}
\newcommand{\cM}{\mathcal{M}}
\newcommand{\cV}{\mathcal{V}}
\newcommand{\cW}{\mathcal{W}}

\newcommand{\CC}{\mathbb{C}}
\newcommand{\NN}{\mathbb{N}}
\newcommand{\RR}{\mathbb{R}}
\newcommand{\ZZ}{\mathbb{Z}}

\renewcommand{\Im}{\mathrm{Im}}
\renewcommand{\Re}{\mathrm{Re}}

\newcommand{\eps}{\epsilon}
\newcommand{\grad}{\nabla}
\newcommand{\lam}{\lambda}
\newcommand{\lap}{\triangle}

\newcommand{\ip}[2]{\ensuremath{\left<#1,#2\right>}}

\newcommand{\image}{\operatorname{im}}
\newcommand{\onull}{\operatorname{null}}
\newcommand{\rank}{\operatorname{rank}}
\newcommand{\range}{\operatorname{range}}
\newcommand{\trace}{\operatorname{tr}}
\newcommand{\Span}{\operatorname{span}}

\newcommand{\prob}[1]{\bigskip\noindent\textbf{#1.}\quad }

\newcommand{\pts}[1]{(\emph{#1 pts}) }
\newcommand{\epart}[1]{\medskip\noindent\textbf{(#1)}\quad }
\newcommand{\ppart}[1]{\,\textbf{(#1)}\quad }

\newcommand{\ds}{\displaystyle}

\newcommand{\nex}{\medskip\noindent}


\begin{document}
\scriptsize \noindent Math 617 Functional Analysis (Bueler) \hfill \emph{assigned 16 February 2026; {\color{red} version 4}}
\normalsize\medskip

\Large\centerline{\textbf{Assignment 4}}
\large
\medskip

\centerline{\textbf{Due Friday 27 February at the beginning of class, \underline{firm}}}
\medskip
\normalsize

\thispagestyle{empty}

\bigskip
\noindent This Assignment is based on Chapters 2, 3, and 4 of our textbook,\footnote{K.~Saxe,~\emph{Beginning Functional Analysis}, Springer 2010.} and on the lectures and slides.\footnote{See the \href{https://bueler.github.io/fa/daily.html}{Daily Log} tab of the public site \href{https://bueler.github.io/fa/index.html}{\texttt{bueler.github.io/fa/}} for the slides.}


\newcommand{\augment}[1]{\quad\,\strut \begin{minipage}[t]{0.7\textwidth} \emph{#1}\end{minipage}}
\newcommand{\augpart}[1]{\emph{\textbf{(#1)}}}

\bigskip
\noindent \textsc{Do the following} Exercises from Chapter 2 (see pages 29--31):

\begin{itemize}
\item \textbf{Exercise 2.2.3}
\end{itemize}

\bigskip
\noindent \textsc{Do the following} Exercises from Chapter 3 (see pages 69--73):

\begin{itemize}
\item \textbf{Exercise 3.2.1}
\item \textbf{Exercise 3.2.9}  \augment{Start by observing that open intervals $O=(a_1,b_1)\times \dots \times (a_n,b_n)$ are already in $\cE$, thus they are in $\cM$.  Then use the fact that $\cM$ is a $\sigma$-ring, i.e.~Theorem 3.6, to prove the things you want.}
\item \textbf{Exercise 3.3.1}
\item \textbf{Exercise 3.3.2}  \augment{You may assume the existence of a non-measurable real set.  Specifically, you may assume that there exists $E\subset[0,1]$ such that $E\notin\cM$.  Now this problem is easy and short.}
\end{itemize}

\bigskip
\noindent \textsc{Do the following} Exercises from Chapter 4 (see pages 88--91):

\begin{itemize}
\item \textbf{Exercise 4.1.5}
\item \textbf{Exercise 4.1.6}
\end{itemize}


\medskip
\noindent \textsc{Do the following additional problems.}

\prob{P5}  On Assignment 3 you did Exercise 4.1.3 to find the classical Fourier series of $f(x)=x^2$ on the interval $[-\pi,\pi]$.  It is a cosine series because the function is even.  Using that result, apply Parseval's equality to show that
    $$\sum_{n=1}^\infty \frac{1}{n^4} = \frac{\pi^4}{90}.$$

\prob{P6}  \ppart{a} Consider a complex sequence $(a_n)_{n=1}^\infty$.  Show that if $a_n \to a \in \CC$ then the new sequence
	$$b_n = \frac{1}{n} \sum_{k=1}^n a_k,$$
which is called the sequence of \emph{Cesaro means}, also converges to $a$.

\epart{b}  Give an example of a real or complex sequence $(a_n)$ which does not converge---justify this---but for which the sequence of Cesaro means $(b_n)$ does converge.

\epart{c}  \emph{This part is a historical research and writing assignment.  You do not need to state the proofs, and no credit will be given for doing so.  Instead, explain the statement of the theorem informally, and explain its significance relative to the theory of Fourier series of continuous functions covered in lectures.  Recall that Dirichlet's 1829 result is stated in my slides for weeks 4 \& 5.}

\medskip \noindent At a young age the Hungarian mathematician Lip\'ot Fej\'er proved \emph{Fej\'er's Theorem}\footnote{L.~Fej\'er, 1910. \emph{Sur le sommes partielles de la serie de Fourier}, Comptes rendus, 150, 1299--1302.} which shows that the coefficients of the Fourier series of a continuous function $f$ can be used to construct a sequence of trigonmetric polynomials which converge uniformly to $f$.  Please look up this theorem, on the internet or in other books, and understand it.  Then explain in complete sentences, at least 5 sentences, but at most 10, what this theorem says, how it uses Cesaro means, and how it is related to what Dirichlet proved.

\end{document}
